\section{Nội dung Đề tài}
    IoT Gateway là một trong những thành phần cốt lõi của hầu hết các hệ thống IoT hiện nay, đóng vai trò cầu nối giữa các giao thức và chuẩn giao tiếp khác nhau. Mỗi ứng dụng cụ thể lại đòi hỏi tập hợp giao thức và chuẩn giao tiếp riêng; vì vậy, mỗi kịch bản sử dụng khác nhau yêu cầu IoT Gateway hỗ trợ các giao thức và chuẩn giao tiếp tương ứng.\par

    Tuy nhiên, việc thiết kế lại từ đầu IoT Gateway cho từng trường hợp riêng lẻ rất tốn kém về thời gian và nhân lực. Để giảm chi phí này, nhiều thiết bị trên thị trường được phát triển theo hướng tích hợp sẵn một dải rộng các giao thức và chuẩn giao tiếp phổ biến. Cách tiếp cận đó giúp đáp ứng nhu cầu đa dạng, nhưng việc duy trì các giao thức và chuẩn giao tiếp không sử dụng lại gây lãng phí và tăng chi phí triển khai.\par
    
    Nhằm giải quyết hạn chế trên, đồ án đặt mục tiêu khảo sát, thiết kế, và hiện thực một IoT Gateway có khả năng cấu hình linh hoạt, tập trung vào hai đặc tính:\par
    \begin{enumerate}
        \item \textbf{Modular} – Phần cứng thiết kế dạng lắp ghép, chỉ gắn các module cần thiết, qua đó loại bỏ phần dư thừa và tối ưu chi phí.
        \item \textbf{Configurability} – Quy trình cấu hình tối giản, thân thiện với người dùng, giúp rút ngắn thời gian triển khai và giảm yêu cầu kỹ năng chuyên sâu.
    \end{enumerate}
    
    Nhờ kết hợp hai đặc tính này, mô hình IoT Gateway đề xuất không chỉ đáp ứng nhu cầu linh hoạt của các hệ thống IoT mà còn tối ưu hóa tổng chi phí sở hữu và vận hành.
